\documentclass[12pt,a4paper]{article}

% Encoding & language
\usepackage[utf8]{inputenc}
\usepackage[american]{babel}

% Typography & layout
\usepackage{geometry}
\geometry{a4paper,left=25mm,right=25mm,top=25mm,bottom=25mm}
\usepackage{setspace}
\doublespacing
\usepackage{microtype}
\usepackage{newtxtext,newtxmath} % Times-like font (APA-friendly)

% Math, graphics, quotes
\usepackage{amsmath}
\usepackage{graphicx}
\usepackage{csquotes}

% References (APA via biblatex + biber)
\usepackage[style=apa,backend=biber,sortcites=true,sorting=nyt]{biblatex}
\DeclareLanguageMapping{american}{american-apa}
\addbibresource{references.bib}

% Links
\usepackage[hidelinks]{hyperref}

% Header/footer: page number top-right (APA student paper)
\usepackage{fancyhdr}
\pagestyle{fancy}
\fancyhf{}
\rhead{\thepage}
\renewcommand{\headrulewidth}{0pt}
\renewcommand{\footrulewidth}{0pt}

% Title metadata
\newcommand{\papertitle}{Security Action for Europe (SAFE):\\
Strengthening EU Strategic Autonomy Through Defense Industrial Integration}
\newcommand{\authorsblock}{Sem Broere (s3410552)\\Nolt Noordhoek\\Coen van der Giesen}
\newcommand{\courseinfo}{European Union in Global Affairs\\Assignment 1}
\newcommand{\groupinfo}{Group 16}
\newcommand{\duedate}{September 29, 2025}

\begin{document}

% Title page
\begin{titlepage}
  \centering
  \vspace*{3cm}
  {\bfseries\Large \papertitle\par}
  \vspace{2cm}
  {\large \authorsblock\par}
  \vspace{1cm}
  {\large \groupinfo\par}
  \vspace{1.5cm}
  {\large \courseinfo\par}
  \vspace{0.75cm}
  {\large Deadline: \duedate\par}
  \vfill
  {\large \today\par}
\end{titlepage}

\setcounter{page}{1}

% Body
\section{Introduction}

The Security Action for Europe (SAFE) represents a watershed moment in European Union defense policy, establishing a €150 billion loan instrument to accelerate defense procurement and industrial capacity. Adopted on May 27, 2025, through Council Regulation (EU) 2025/1106, SAFE emerged as the EU's response to deteriorating security conditions following Russia's full-scale invasion of Ukraine and evolving global threats \parencite{council2025safe, commission2025proposal}. This paper examines SAFE's decision-making process, institutional dynamics, and contribution to EU power projection in the international arena. Through analysis of its implementation mechanisms and emerging challenges, this research demonstrates how SAFE strengthens the EU's strategic autonomy while revealing persistent tensions between emergency decision-making and democratic legitimacy in European defense integration.

\section{The SAFE Instrument: Objectives and Framework}

SAFE constitutes the EU's most ambitious defense financing mechanism to date, providing long-term, competitively priced loans to Member States for defense procurement with 45-year maximum durations and 10-year grace periods \parencite{commission2025proposal}. The instrument addresses critical capability gaps identified in the 2025 European Defense Industrial Strategy, particularly ammunition stockpiles, air defense systems, and strategic enablers. As outlined in COM(2025) 122 final, SAFE operates through a two-tier categorization system: Category 1 encompasses immediate needs including ammunition and cyber capabilities, while Category 2 covers complex systems like air defense and maritime platforms requiring enhanced European industrial participation \parencite{commission2025safe}.

The measure directly links to the broader ReArm Europe Plan, an €800 billion defense spending initiative through 2030. SAFE represents the first operational pillar, complemented by Stability and Growth Pact flexibility measures, cohesion fund adaptations, and private capital mobilization mechanisms \parencite{euronews2025}. This comprehensive approach signals the EU's shift from fragmented national procurement toward coordinated European defense industrial policy, addressing long-standing inefficiencies where Member States collectively spend four times more than Russia on defense yet achieve significantly lower output in critical capabilities.

Central to SAFE's design is the European preference requirement mandating minimum 65\% content from EU, EEA-EFTA countries, or Ukraine. This provision aims to strengthen the European Defense Technological and Industrial Base (EDTIB) while reducing dependencies on third-country suppliers \parencite{commission2025proposal}. The regulation includes VAT exemptions for SAFE-supported procurements and enables 15\% pre-financing upon agreement signature, addressing immediate liquidity needs for rapid capability development.

\section{Decision-Making Process and Institutional Dynamics}

The adoption of SAFE revealed fundamental tensions in EU crisis decision-making. The Council employed Article 122 TFEU, the Treaty's solidarity clause for \enquote{severe difficulties caused by natural disasters or exceptional occurrences beyond control.} This emergency procedure enabled unanimous Council adoption without European Parliament co-legislation, marking only the second use of Article 122 for defense matters following the 2022 ammunition procurement initiative \parencite{verfassungsblog2025}.

The Commission's proposal on March 19, 2025, underwent accelerated Council examination, with the Hungarian Presidency prioritizing rapid adoption. All 27 Member States supported the regulation during the May 27 Council meeting, demonstrating remarkable unanimity given traditional divisions on defense matters \parencite{council2025safe}. However, this consensus masked underlying procedural concerns. Several Member States, while supporting SAFE's objectives, expressed reservations about circumventing ordinary legislative procedures, particularly given the instrument's scale and multi-year implications.

The European Parliament's exclusion from co-decision triggered immediate institutional conflict. On August 25, 2025, Parliament
\begin{document}